\chapter{Geodezyjne, odwzorowanie wykładnicze i zupełność}
\label{ch:geodezyjne-exp}

W Rozdziale~\ref{ch:koneksje} otrzymaliśmy równanie
$D\dot\gamma/\dd t=0$ jako warunek stacjonarności
energii. Teraz wyjaśnimy, kiedy taka krzywa
rzeczywiście jest najkrótsza, jak opisuje otoczenie
punktu i co gwarantuje jej istnienie przez dowolnie
długi czas. Ostatnia część rozdziału odróżnia te
riemannowskie twierdzenia od geodezyjnych
czasoprzestrzeni Schwarzschilda i Kerra.

W całym rozdziale rozmaitości są gładkie i \emph{bez brzegu}.
Jeśli nie zaznaczono inaczej, metryka jest riemannowska;
odległości rozpatrujemy w jednej składowej spójnej.
Brak brzegu jest istotny: domknięty przedział z metryką
euklidesową jest zwarty, ale prosta dochodząca do jego
końca nie ma geodezyjnego przedłużenia w tym przedziale.

\section{Geodezyjna i pierwsza wariacja}
\label{sec:geodezyjna-pierwsza-wariacja}

\begin{definicja}[Geodezyjna]
\index{geodezyjna@Geodezyjna}
\label{def:geodezyjna}
Geodezyjną koneksji $\nabla$ nazywamy krzywą
$\gamma:I\to M$ spełniającą
\begin{equation}\label{eq:geodezyjna-ode}
 \nabla_{\dot\gamma}\dot\gamma=0,
 \qquad
 \ddot x^k+\Gamma^k_{ij}(x)\dot x^i\dot x^j=0.
\end{equation}
Użyty parametr jest \emph{afiniczny}.
Gdy $g$ jest metryką, mamy na myśli jej
koneksję Leviego-Civity.
\end{definicja}

Zgodność koneksji z metryką daje
$\frac{\dd}{\dd t}g(\dot\gamma,\dot\gamma)
=2g(D\dot\gamma/\dd t,\dot\gamma)=0$.
Wektor prędkości geodezyjnej riemannowskiej ma więc stałą
normę. Parametr $s=at+b$, $a\ne0$,
zachowuje równanie~\eqref{eq:geodezyjna-ode};
dowolna nieliniowa zmiana parametru na ogół
go nie zachowuje. Dla przykładu prosta
$\gamma(t)=p+tv$ w $\R^n$ jest geodezyjną,
ale dla $v\ne0$ krzywa $p+t^2v$ nie spełnia równania geodezyjnej
w parametrze $t$.

\begin{stwierdzenie}[Dane początkowe]
\label{prop:geodezyjna-dane-poczatkowe}
Dla każdego $p\in M$ i $v\in T_pM$ istnieje
jedyna maksymalna geodezyjna $\gamma_{p,v}$
o $\gamma_{p,v}(0)=p$, $\dot\gamma_{p,v}(0)=v$.
Jest określona co najmniej dla $|t|<\varepsilon$
zależnego od $(p,v)$ i zależy gładko od danych
tam, gdzie jest określona.
\end{stwierdzenie}
\begin{proof}
W mapie wprowadzamy $u^k=\dot x^k$.
Równanie drugiego rzędu staje się gładkim
układem pierwszego rzędu na otwartym
podzbiorze $\R^{2n}$:
\[
 \dot x^k=u^k,
 \qquad \dot u^k=-\Gamma^k_{ij}(x)u^iu^j.
\]
Twierdzenie Picarda--Lindelöfa o lokalnym
istnieniu, jednoznaczności i gładkiej zależności
od danych dla gładkiego pola na $TM$ daje
rozwiązanie. Jednoznaczność pozwala skleić
rozwiązania z różnych map i przedłużać je
do maksymalnego przedziału. Nie zakładamy
jeszcze, że tym przedziałem jest cała $\R$.
\end{proof}

Pierwsza wariacja jest dobrym miejscem na
rozróżnienie dwóch zdań: \emph{geodezyjna jest
punktem stacjonarnym długości} oraz
\emph{krótki jej fragment minimalizuje długość}.
Drugie wymaga osobnego dowodu poniżej.

\begin{twierdzenie}[Wzór na pierwszą wariację długości]
\label{thm:pierwsza-wariacja-dlugosci}
Niech $F(s,t)$ będzie gładką wariacją regularnej
krzywej $\gamma(t)=F(0,t)$, $t\in[a,b]$, $a<b$,
na rozmaitości riemannowskiej.
Oznaczmy $V=\partial_sF|_{s=0}$ i
$T=\dot\gamma/|\dot\gamma|_g$.
Wtedy
\begin{equation}\label{eq:pierwsza-wariacja-dlugosci}
 \left.\frac{\dd}{\dd s}L(F_s)\right|_{s=0}
  =\bigl[g(V,T)\bigr]_a^b
    -\int_a^b g\!\left(V,\frac{DT}{\dd t}\right)\dd t.
\end{equation}
Jeżeli $\gamma$ ma prędkość $1$, to
$T=\dot\gamma$ i $DT/\dd t=D\dot\gamma/\dd t$.
Przy ustalonych końcach składnik brzegowy znika.
\end{twierdzenie}
\begin{proof}
Połóżmy $W=\partial_tF$. Regularność $\gamma$ i zwartość
$[a,b]$ dają dodatnie dolne ograniczenie $|W(0,t)|_g$.
Po zmniejszeniu zakresu $s$ także $|W(s,t)|_g$ jest
jednostajnie dodatnia, więc możemy różniczkować pierwiastek.
Różniczkujemy
pod całką $L(F_s)=\int_a^b\sqrt{g(W,W)}\,\dd t$.
Zgodność koneksji z metryką daje
\[
 \partial_s|W|_g
 =\frac{g(D_sW,W)}{|W|_g}
 =g(D_sW,T).
\]
Pochodne $D_s,D_t$ bierzemy wzdłuż odpowiednich krzywych
w rodzinie $F$. Nie trzeba zakładać, że $F$ jest immersją.
Współrzędne dają bezpośrednio
\[
 (D_s\partial_tF)^k-(D_t\partial_sF)^k
 =(\Gamma^k_{ij}-\Gamma^k_{ji})
                      \partial_sF^i\partial_tF^j=0,
\]
ponieważ pochodne mieszane komutują, a koneksja nie ma
torsji. W $s=0$ otrzymujemy $D_sW=D_tV$.
Z reguły iloczynu
\[
 g(D_tV,T)=\frac{\dd}{\dd t}g(V,T)-g(V,D_tT).
\]
Całkujemy tę równość i podstawiamy do
pochodnej długości. To daje dokładnie
\eqref{eq:pierwsza-wariacja-dlugosci}.
W ciągłej wariacji kawałkami gładkiej, z ustalonym
skończonym podziałem i regularnymi kawałkami, stosujemy
wzór na każdym kawałku. W punkcie sklejenia $t_0$
pozostaje dodatkowo
$g(V(t_0),T(t_0^-)-T(t_0^+))$:
wyjaśnimy poniżej, jak ten wyraz pozwala skrócić narożnik.
\end{proof}

\begin{wniosek}[Dlaczego krzywa minimalna musi być geodezyjną]
\label{cor:minimalna-geodezyjna}
Regularna, gładka krzywa lokalnie minimalizująca
długość między swoimi końcami jest geodezyjną
po parametryzacji długością łuku.
\end{wniosek}
\begin{proof}
Parametryzujmy ją ze stałą prędkością $1$.
Dla wariacji o ustalonych końcach pochodna
długości jest zerowa. W mapie możemy wybrać
pole wariacyjne $V$ o zwartym nośniku
w dowolnie małym odcinku i dowolnych
składowych: realizuje je na przykład
$F(s,t)=x^{-1}(x(\gamma(t))+sV^i(t)e_i)$
dla małego $s$. Wzór~\eqref{eq:pierwsza-wariacja-dlugosci}
i lemat podstawowy rachunku wariacyjnego
dają $DT/\dd t=0$ w każdym punkcie.
\end{proof}

\section{Odwzorowanie wykładnicze i współrzędne normalne}
\label{sec:exp-normalne}

Zaczynamy w jednym punkcie $p$ i wystrzeliwujemy
geodezyjne we wszystkich kierunkach $v\in T_pM$.
Wektor zawiera \emph{kierunek i skalę}: geodezyjna
startująca z $v$ przebywa w czasie $1$ drogę
o długości $|v|_g$, o ile jest określona.

\begin{definicja}[Odwzorowanie wykładnicze]
\index{odwzorowanie wykladnicze@Odwzorowanie wykładnicze}
\label{def:exp-geodezyjna}
Niech $D_p\subset T_pM$ będzie zbiorem wektorów
$v$, dla których $\gamma_{p,v}$ jest określona
na całym $[0,1]$. Kładziemy
\begin{equation}\label{eq:exp-definicja}
 \exp_p:D_p\longrightarrow M,
 \qquad \exp_p(v)=\gamma_{p,v}(1).
\end{equation}
To odwzorowanie nie jest na ogół globalnie
określone ani różnowartościowe. Zbiór $D_p$ jest otwarty,
zawiera $0$ i jest gwiaździsty względem $0$. Istotnie,
maksymalny przepływ gładkiego pola geodezyjnego na $TM$
ma otwartą dziedzinę w $\R\times TM$; warunek istnienia
w chwili $1$ jest zatem otwarty w danych początkowych.
Istnienie od $0$ do $1$ wynika z tego, że przedział
maksymalnego rozwiązania jest przedziałem zawierającym $0$.
Krzywa stała daje $0\in D_p$.

Jednoznaczność i podstawienie do równania dają prawo
$\gamma_{p,cv}(t)=\gamma_{p,v}(ct)$ dla dowolnej stałej $c$
na odpowiadających sobie przedziałach istnienia.
Dla $v\in D_p$ i $0\leq t\leq1$ wynika stąd
gwiaździstość oraz
\begin{equation}\label{eq:exp-skalowanie}
 \exp_p(tv)=\gamma_{p,v}(t).
\end{equation}
Gładkość przepływu daje gładkość $\exp_p$ na całym $D_p$,
a także gładką zależność od $p$ na otwartym podzbiorze $TM$.
\end{definicja}

\begin{twierdzenie}[Lokalne istnienie współrzędnych normalnych]
\label{thm:wspolrzedne-normalne}
Dla każdego $p\in(M,g)$ istnieje $\rho>0$
takie, że $B_\rho(0)\subset T_pM$ leży w $D_p$
i $\exp_p$ ograniczone do tego zbioru jest
dyfeomorfizmem na otoczenie $U$ punktu $p$.
Po wybraniu ortonormalnej bazy $T_pM$
odwrotność $\exp_p^{-1}:U\to\R^n$ jest
\emph{mapą we współrzędnych normalnych}.
W jej środku
\begin{equation}\label{eq:normalne-metryka}
 g_{ij}(p)=\delta_{ij},\qquad
 \Gamma^k_{ij}(p)=0,\qquad
 \partial_\ell g_{ij}(p)=0.
\end{equation}
Stwierdzenie lokalnego dyfeomorfizmu działa
również dla metryki pseudoriemannowskiej;
w bazie pseudoortonormalnej $g_{ij}(p)$ zastępuje wtedy
macierz diagonalna o odpowiedniej sygnaturze. Kulę
w przestrzeni stycznej wybieramy za pomocą pomocniczej
normy euklidesowej, ponieważ metryka nieokreślona nie jest normą.
\end{twierdzenie}
\begin{proof}
Gładka zależność geodezyjnych od danych
początkowych w Stwierdzeniu~\ref{prop:geodezyjna-dane-poczatkowe}
gwarantuje, że dla $v$ bliskich $0$ wszystkie
rozwiązania istnieją aż do czasu $1$.
Istotnie stała krzywa $\gamma_{p,0}(t)=p$
istnieje dla każdego $t$; na zwartym odcinku
$[0,1]$ można wybrać wspólne małe otoczenie
jej danych, w którym działa twierdzenie o gładkiej
zależności ODE. Wtedy $\exp_p$ jest gładkie.

Niech $w\in T_pM$. Z \eqref{eq:exp-skalowanie}
wynika $\exp_p(sw)=\gamma_{p,w}(s)$ dla małego $s$.
Różniczkując w $s=0$, dostajemy
$\dd(\exp_p)_0(w)=\dot\gamma_{p,w}(0)=w$.
Zatem $\dd(\exp_p)_0$ jest identycznością
$T_0(T_pM)\simeq T_pM$. Twierdzenie o funkcji
odwrotnej daje mniejszą kulę, na której
$\exp_p$ jest dyfeomorfizmem.

W ortonormalnej bazie jej różniczka w $0$
identyfikuje bazę współrzędnych z wybraną bazą
$T_pM$, więc $g_{ij}(p)=\delta_{ij}$.
Krzywe $x(t)=tv$ są geodezyjnymi w tej mapie.
Wstawiając je do \eqref{eq:geodezyjna-ode}
w $t=0$, dostajemy
$\Gamma^k_{ij}(p)v^iv^j=0$ dla każdego $v$.
Ponieważ koneksja Leviego-Civity nie ma torsji,
$\Gamma^k_{ij}=\Gamma^k_{ji}$; z polaryzacji
tej równości wynika $\Gamma^k_{ij}(p)=0$.
Wreszcie zgodność z metryką mówi
$\partial_\ell g_{ij}
=g_{mj}\Gamma^m_{\ell i}
 +g_{im}\Gamma^m_{\ell j}$, zatem pochodne
metryki znikają w $p$.
\end{proof}

\begin{figure}[htbp]
\centering
\begin{tikzpicture}[>=Latex,scale=.95]
  \fill[manifoldgold!8] (-3.3,0) circle (1.34);
  \draw[manifoldgold!65!black,thick] (-3.3,0) circle (1.34);
  \fill[manifoldblue!8] (2,0) ellipse (2.15 and 1.28);
  \draw[manifoldblue!65!black,thick] (2,0) ellipse (2.15 and 1.28);
  \coordinate (z) at (-3.3,0);
  \coordinate (pa) at (1.35,-.24);
  \fill[black] (z) circle (2pt) node[below left] {$0$};
  \fill[black] (pa) circle (2pt) node[below left] {$p$};
  \draw[manifoldred,very thick,->] (z)--(-2.34,.63) node[above right] {$v$};
  \draw[manifoldgreen!75!black,very thick,->] (z)--(-3.85,.86) node[above left] {$w$};
  \draw[manifoldred,very thick,->] (pa) .. controls (1.84,-.24) and (2.54,.26) .. (3.02,.61);
  \draw[manifoldgreen!75!black,very thick,->] (pa) .. controls (1.10,.23) and (.94,.56) .. (1.08,.88);
  \fill[manifoldred] (3.02,.61) circle (2pt) node[above right] {$\exp_p(v)$};
  \fill[manifoldgreen!75!black] (1.08,.88) circle (2pt) node[above left] {$\exp_p(w)$};
  \draw[manifoldblue!80!black,very thick,->] (-1.52,0)--(-.62,0)
       node[midway,above] {$\exp_p$};
  \node[manifoldgold!70!black] at (-3.3,-1.68) {$T_pM$};
  \node[manifoldblue!80!black] at (2,-1.68) {$U\subset M$};
\end{tikzpicture}
\caption{Wektory w $T_pM$ opisują krótkie geodezyjne wychodzące z $p$. Odwzorowanie $\exp_p$ zamienia małą kulę wektorów na otoczenie punktu.}
\label{fig:exp-normalne}
\end{figure}

Słowo „normalne” dotyczy tutaj współrzędnych,
nie wektora normalnego do powierzchni.
Równość $\Gamma(p)=0$ zachodzi w \emph{jednym}
punkcie. Nie daje podstaw do wyzerowania
wszystkich symboli na całej zakrzywionej
rozmaitości.

\section{Lemat Gaussa i lokalna minimalność}
\label{sec:gauss-geodezyjna-minimalna}

Współrzędne normalne pozwalają rozdzielić
ruch \emph{od} $p$ od ruchu \emph{w poprzek}.
Lemat Gaussa mówi, że te kierunki są
ortogonalne w metryce, a nie tylko na rysunku.

\begin{lemat}[Gaussa]
\label{lem:gaussa-exp}
Jeżeli $v$ leży w otwartej dziedzinie $D_p$ odwzorowania
$\exp_p$ i
$w\in T_pM$, to
\begin{equation}\label{eq:lemat-gaussa}
 g_{\exp_p(v)}\!\left(
   \dd(\exp_p)_v(v),\dd(\exp_p)_v(w)\right)
   =g_p(v,w).
\end{equation}
W szczególności pochodna w kierunku radialnym
jest prostopadła do obrazów kierunków
stycznych do sfery $|v|_g=\mathrm{const}$.
\end{lemat}
\begin{proof}
Rozpatrzmy dwuwymiarową rodzinę geodezyjnych
$F(s,t)=\exp_p(t(v+sw))$, $0\leq t\leq1$,
przy małym $s$. Połóżmy
$T=\partial_tF$, $V=\partial_sF$.
Każda krzywa $t\mapsto F(s,t)$ jest geodezyjną,
więc $D_tT=0$, a jej kwadrat prędkości
jest stały i równy $g_p(v+sw,v+sw)$.
Z braku torsji $D_tV=D_sT$. Zatem
\[
 \partial_t g(V,T)
 =g(D_sT,T)+g(V,D_tT)
 =\tfrac12\partial_s g(T,T).
\]
W $s=0$ prawa strona jest stałą $g_p(v,w)$.
Ponadto $V(0)=0$, gdyż $F(s,0)=p$.
Całkując po $t$ od $0$ do $1$, dostajemy
$g(V(1),T(1))=g_p(v,w)$.
Reguła łańcuchowa daje
$V(1)=\dd(\exp_p)_v(w)$ i
$T(1)=\dd(\exp_p)_v(v)$, czyli tezę.
\end{proof}

W normalnym otoczeniu zapiszmy niezerowe
$v=r\theta$, gdzie $r=|v|_g$ i
$|\theta|_g=1$. Lemat daje lokalnie postać
\begin{equation}\label{eq:metryka-biegunowo-normalna}
 g=\dd r^2+h_{ab}(r,\theta)
                 \dd\theta^a\dd\theta^b.
\end{equation}
Rzeczywiście $\partial_r=\dd\exp_{r\theta}(\theta)$
ma długość $1$ (stała prędkość radialnej
geodezyjnej), a kierunki kątowe są do niego
prostopadłe. Macierz $(h_{ab})$ jest dodatnio
określona dla $r>0$ w tym otoczeniu.

\begin{twierdzenie}[Krótka geodezyjna minimalizuje]
\label{thm:lokalna-minimalnosc-geodezyjnej}
Istnieje $\rho>0$ takie, że dla każdego
$q=\exp_p(v)$ z $|v|_g<\rho$ radialna
geodezyjna $\gamma(t)=\exp_p(tv)$,
$0\leq t\leq1$, spełnia
\begin{equation}\label{eq:odleglosc-exp-lokalna}
 d_g(p,q)=L_g(\gamma)=|v|_g.
\end{equation}
Dla $v\ne0$ każda odcinkami $C^1$ droga osiągająca tę
długość ma ten sam ślad radialny i przebiega go bez cofania;
może przy tym zatrzymywać się na pewnych odcinkach parametru.
Dla $v=0$ jedyną drogą o długości zero jest droga stała. W szczególności \emph{każdy}
dostatecznie krótki fragment dowolnej
geodezyjnej riemannowskiej minimalizuje
odległość między swoimi końcami.
\end{twierdzenie}
\begin{proof}
Wybierzmy $\rho$ tak małe, by
$\exp_p$ było dyfeomorfizmem na kuli $B_{3\rho}(0)$.
W szczególności domknięcie $\overline B_{2\rho}(0)$
leży w tej dziedzinie, a jego obraz jest zwarty. W normalnych współrzędnych
dowolną odcinkami $C^1$ krzywą $\alpha$
biegnącą od $p$ do $q$ i pozostającą
w obrazie tej kuli piszemy dla $r>0$
przez $r(t),\theta(t)$. Z \eqref{eq:metryka-biegunowo-normalna}
wynika
\[
 |\dot\alpha|_g^2
 =(\dot r)^2+h_{ab}\dot\theta^a\dot\theta^b
 \geq (\dot r)^2,
 \qquad
 L_g(\alpha)\geq\int|\dot r|\,\dd t
 \geq r(q)=|v|_g.
\]
W punkcie $p$ funkcja $r$ nie musi być gładka.
Jednak $u(t)=\exp_p^{-1}(\alpha(t))$ jest odcinkami $C^1$,
a $r(t)=|u(t)|_g$ jest lipschitzowska na zwartym przedziale,
więc absolutnie ciągła. Pochodna $\dot r$ istnieje prawie
wszędzie i jest równa zero prawie wszędzie na zbiorze
$\{t:r(t)=0\}$: w każdym wewnętrznym punkcie tego zbioru,
w którym pochodna istnieje, $r$ ma minimum. Nierówność
$|\dot\alpha|_g\geq|\dot r|$ zachodzi zatem prawie wszędzie,
co uzasadnia całkowanie również przy wielokrotnych powrotach do $p$. Radialna geodezyjna
ma $\dot\theta=0$, $\dot r\geq0$ i długość
$|v|_g$, więc osiąga dolne ograniczenie.
Równość w drugiej nierówności całkowej wymaga
$\dot r\geq0$ prawie wszędzie, czyli $r$ jest niemalejąca.
Zbiór $\{r>0\}$ jest wtedy jednym końcowym przedziałem.
Na nim równość w pierwszej nierówności i dodatnia
określoność $h$ dają $\dot\theta=0$ prawie wszędzie.
Na każdym zwartym podprzedziale $\theta$ jest odcinkami
$C^1$, a więc stała. To dowodzi jednoznaczności śladu;
chwile postoju odpowiadają stałym odcinkom funkcji $r$.

Jeśli $\alpha$ opuści obraz $B_{2\rho}(0)$,
jej pierwszy punkt wyjścia ma promień $2\rho$.
Poprzednia nierówność na początkowym kawałku
daje $L_g(\alpha)\geq2\rho>|v|_g$.
Porównaliśmy więc także drogi wychodzące
z otoczenia, co jest potrzebne dla odległości
$d_g$ zdefiniowanej przez \emph{wszystkie}
dopuszczalne krzywe. Dla dowolnej geodezyjnej
powtarzamy argument, biorąc za $p$ jej punkt
początkowy i skracając przedział parametru.
\end{proof}

Z dowodu wynika również identyfikacja małych kul. Dla
$0<a<\rho$ mamy
\begin{equation}\label{eq:male-kule-normalne}
 B^g_a(p)=\exp_p(B_a(0)),\qquad
 \overline B^g_a(p)=\exp_p(\overline B_a(0)).
\end{equation}
Tutaj $\overline B^g_a(p)=\{q:d_g(p,q)\leq a\}$.
Każdy punkt spoza $\exp_p(B_{2\rho}(0))$ ma odległość
co najmniej $2\rho$ od $p$, na mocy argumentu z pierwszym
wyjściem. Wewnątrz tego otoczenia ten sam rachunek
radialny daje $d_g(p,\exp_p(v))=|v|_g$.
Stąd obie równości i zwartość małych kul domkniętych
oraz sfer metrycznych, bez założenia zupełności $M$.

Pierwsza wariacja dowodziła, że krzywa minimalna
\emph{musi} być geodezyjną; lemat Gaussa dowiódł,
że krótka geodezyjna \emph{jest} krzywą minimalną.
Na długich odcinkach drugie zdanie może
przestać być prawdziwe. Druga wariacja i pola
Jacobiego w Sekcji~\ref{sec:druga-wariacja-jacobi}
opiszą, jak krzywizna wpływa na utratę tej własności.

\section{Geodezyjne w geometriach modelowych}
\label{sec:geodezyjne-modele}

\begin{przyklad}[Przestrzeń euklidesowa]
\label{prz:geodezyjne-euklidesowe}
Dla $g=\sum_i(\dd x^i)^2$ symbole Christoffela
znikają. Równanie $\ddot x=0$ ma rozwiązanie
$\gamma_{p,v}(t)=p+tv$, a $\exp_p(v)=p+v$
na całym $T_p\R^n$. Dla dowolnej krzywej
$\alpha:[a,b]\to\R^n$ od $p$ do $q$
nierówność trójkąta dla całki daje
\[
 L(\alpha)=\int_a^b|\dot\alpha|\,\dd t
 \geq\left|\int_a^b\dot\alpha\,\dd t\right|
 =|q-p|.
\]
Odcinek prosty osiąga tę wartość i jest
globalnie minimalny.
\end{przyklad}

\begin{figure}[htbp]
\centering
\begin{tikzpicture}[>=Latex,scale=1]
  \fill[manifoldblue!5,rounded corners=8pt]
     (-3.4,-1.05) rectangle (3.4,1.05);
  \draw[manifoldblue!45,rounded corners=8pt]
     (-3.4,-1.05) rectangle (3.4,1.05);
  \coordinate (p) at (-2.45,-.48);
  \coordinate (q) at (2.46,.51);
  \draw[manifoldred!85!black,very thick,->]
     (p)--(q) node[pos=.50,below] {$\gamma(t)=p+tv$};
  \draw[manifoldblue!65,thick,densely dashed]
     (p) .. controls (-1.68,1.11) and (1.22,1.12) .. (q);
  \draw[manifoldblue!65,thick,densely dashed]
     (p) .. controls (-1.2,-1.24) and (1.35,-.88) .. (q);
  \fill[manifoldblue!80!black] (p) circle (2.5pt)
     node[below left] {$p$};
  \fill[manifoldblue!80!black] (q) circle (2.5pt)
     node[above right] {$q$};
\end{tikzpicture}
\caption{Euklidesowa geodezyjna między $p$ i $q$ jest prostym odcinkiem. Przykładowe drogi przerywane mają większą długość.}
\label{fig:geodezyjna-euklidesowa}
\end{figure}

\begin{przyklad}[Sfera: wielkie koła i antypody]
\label{prz:geodezyjne-sfera}
Niech $R>0$, $p\in S_R^2\subset\R^3$ i $v\perp p$,
$s=|v|$. Dla $v\ne0$ geodezyjna i odwzorowanie
wykładnicze mają postać
\begin{equation}\label{eq:exp-sfera}
 \gamma_{p,v}(t)
 =\cos\!\left(\frac{st}{R}\right)p
  +\frac{R}{s}\sin\!\left(\frac{st}{R}\right)v,
 \qquad \exp_p(v)=\gamma_{p,v}(1).
\end{equation}
Dla $v=0$ rozumiemy granicę, czyli $\exp_p(0)=p$.
Różniczkując dwa razy, dostajemy
$\ddot\gamma=-(s^2/R^2)\gamma$.
Wektor $\gamma$ jest normalny do sfery,
więc jego rzut na przestrzeń styczną jest zerowy.
Ze Stwierdzenia~\ref{prop:koneksja-projekcja} wynika
$D\dot\gamma/\dd t=0$. Są to łuki wielkich kół.

Dla bieguna $p$ i $r=R\theta$ metryka
\eqref{eq:metryka-sfery} przyjmuje postać
\[
 g=\dd r^2+R^2\sin^2(r/R)\,\dd\phi^2,
 \qquad 0<r<\pi R.
\]
Jest to konkretny przypadek
\eqref{eq:metryka-biegunowo-normalna}.
Łuk od $p$ do punktu w odległości kątowej
$\theta<\pi$ ma długość $R\theta$ i jest
minimalny. Funkcja
$r(x)=R\arccos(\langle p,x\rangle/R^2)$
jest określona na całej sferze i gładka
poza biegunami $p,-p$. Lokalny wzór na metrykę
działa po zmianie kąta $\phi$ także przy
przekroczeniu szwu mapy; dla drogi omijającej
antypodę nadal $L\geq\int|\dot r|\geq R\theta$.
Dla drogi przechodzącej przez antypodę najpierw stosujemy
to oszacowanie do pierwszego osiągnięcia poziomu
$r=\pi R-\varepsilon$. Przejście $\varepsilon\downarrow0$
daje długość co najmniej $\pi R>R\theta$.
Przy przejściach przez $p$ argument absolutnej ciągłości
z dowodu Twierdzenia~\ref{thm:lokalna-minimalnosc-geodezyjnej}
uzasadnia całkowanie mimo niegładkości $r$.
Przy $r=\pi R$ wszystkie kierunki jednostkowe
trafiają do antypody $-p$: istnieje wiele
minimalnych półokręgów o długości $\pi R$.
Za antypodą przedłużony łuk nadal jest
geodezyjną, lecz nie minimalizuje długości.
Sfera pozostaje geodezyjnie zupełna: wzór
\eqref{eq:exp-sfera} działa dla każdego $t\in\R$.
\end{przyklad}

\begin{figure}[htbp]
\centering
\begin{tikzpicture}[>=Latex,scale=1.04]
  % Rzut (X,Y,Z) -> (X,Z); oba półokręgi mają Y>=0.
  \shade[inner color=white,outer color=manifoldblue!35] (0,0) circle (1.75);
  \draw[manifoldblue!75!black,thick] (0,0) circle (1.75);
  \draw[manifoldblue!45] (-1.75,0)--(1.75,0);
  \draw[manifoldred,ultra thick,->]
    plot[domain=0:180,samples=80,variable=\u]
      ({1.75*sin(40)*sin(\u)},{1.75*cos(\u)});
  \draw[manifoldgreen!75!black,ultra thick,->]
    plot[domain=0:180,samples=80,variable=\u]
      ({-1.75*sin(40)*sin(\u)},{1.75*cos(\u)});
  \fill[manifoldblue!80!black] (0,1.75) circle (2.5pt) node[above] {$p$};
  \fill[manifoldblue!80!black] (0,-1.75) circle (2.5pt) node[below] {$-p$};
  \node[manifoldred!85!black] at (2.02,.55) {$\gamma_1$};
  \node[manifoldgreen!65!black] at (-2.02,.55) {$\gamma_2$};
\end{tikzpicture}
\caption{Dwa różne półokręgi wielkich kół łączą $p$ z antypodą.
Każdy ma długość $\pi R$. Rysunek jest rzutem ortogonalnym:
półokręgi widać jako połówki elips, a równik jako odcinek.
Przedłużenie łuku za antypodę nie minimalizuje już długości od $p$.}
\label{fig:geodezyjne-antypody}
\end{figure}

\begin{przyklad}[Przestrzeń hiperboliczna]
\label{prz:geodezyjne-hiperboliczne}
W modelu górnej półprzestrzeni z
\eqref{eq:hip-polprzestrzen} pionowa krzywa
\[
 \gamma(s)=(x_0,y_0e^s),\qquad s\in\R,
\]
ma prędkość hiperboliczną $1$.
W dwuwymiarowym przypadku wzory
\eqref{eq:gamma-hiperboliczna} dają
$\ddot y-(\dot y)^2/y=0$ i równanie
poziome jest spełnione, więc to geodezyjna.
Dla krzywej łączącej $(x_0,y_0)$ z
$(x_0,y_1)$ mamy
\[
 L=\int\frac{\sqrt{|\dot x|^2+\dot y^2}}{y}\,\dd t
 \geq\int\left|\frac{\dot y}{y}\right|\dd t
 \geq\left|\log\frac{y_1}{y_0}\right|.
\]
Pionowa geodezyjna osiąga tę wartość.
Zatem granica $y=0$ jest w nieskończonej
odległości wzdłuż pionu, choć na rysunku
euklidesowym wydaje się bliska.

To nie są wszystkie geodezyjne półprzestrzeni.
W wymiarze dwa równania geodezyjnej z
\eqref{eq:gamma-hiperboliczna} mają postać
\[
 \ddot x-\frac{2\dot x\dot y}{y}=0,
 \qquad
 \ddot y+\frac{\dot x^2-\dot y^2}{y}=0.
\]
Rozpatrzmy niestałą geodezyjną i parametryzujmy ją długością łuku:
$(\dot x^2+\dot y^2)/y^2=1$. Pierwsze równanie
po przekształceniu daje
$\frac{\dd}{\dd s}(\dot x/y^2)=0$.
Jeżeli stała $c=\dot x/y^2$ jest zerowa,
otrzymujemy pionową prostą. Jeśli $c\ne0$,
to $\dot x=cy^2$ oraz
$\dot y^2=y^2(1-c^2y^2)$. Eliminując parametr,
\[
 \frac{\dd x}{\dd y}
   =\pm\frac{cy}{\sqrt{1-c^2y^2}},
 \qquad
 (x-x_c)^2+y^2=c^{-2}
\]
po scałkowaniu na odcinkach monotoniczności $y$.
Jest to półokrąg o środku $(x_c,0)$,
przecinający brzeg $y=0$ prostopadle.
Punkt zwrotny wysokości nie jest końcem geodezyjnej.
Dla $R_0=1/|c|$ jawne rozwiązanie jednostkowe to
\[
 x(s)=x_c+\operatorname{sgn}(c)R_0\tanh(s-s_0),\qquad
 y(s)=\frac{R_0}{\cosh(s-s_0)},\qquad s\in\R.
\]
Podstawienie sprawdza oba równania i stałą $\dot x/y^2=c$;
rozwiązanie przechodzi gładko przez maksimum $y$ w $s=s_0$.
W wymiarze $n>2$ każda geodezyjna mieści się
w pewnej pionowej dwupłaszczyźnie: odbicie
w niej jest izometrią, zachowuje punkt i wektor
początkowy, a jednoznaczność rozwiązania
równania geodezyjnej wymusza niezmienniczość
krzywej. Tam stosujemy powyższy rachunek. W wymiarze $1$
model jest półprostą z metryką $\dd y^2/y^2$ i pozostają
wyłącznie geodezyjne pionowe oraz krzywe stałe.

W modelu hiperboloidy z
Przykładu~\ref{prz:trzy-modele-hiperboliczne}
piszemy $\langle(t,z),(t',z')\rangle_L
=-tt'+z\cdot z'$ i $\langle p,p\rangle_L=-1$.
Dla $v\in T_p\mathcal H^n$, $s=|v|_g$, mamy
\begin{equation}\label{eq:exp-hiperboloida}
 \exp_p(v)=\cosh(s)p+\frac{\sinh(s)}{s}v,
\end{equation}
z wartością graniczną $p$ przy $s=0$.
Dokładniej, dla jednostkowego $w\in T_p\mathcal H^n$
cała geodezyjna ma wzór
$\gamma_{p,w}(\tau)=\cosh\tau\,p+\sinh\tau\,w$.
Leży w przecięciu hiperboloidy z dwuwymiarową
płaszczyzną liniową $\operatorname{span}(p,w)$
przechodzącą przez początek przestrzeni Minkowskiego.
Rzeczywiście $\langle p,v\rangle_L=0$;
podstawienie daje normę lorentzowską $-1$,
a znak współrzędnej czasowej pozostaje dodatni przez ciągłość.
Rzutowanie pochodnej płaskiej działa tu także dla iloczynu
lorentzowskiego: przestrzeń styczna jest niezdegenerowana,
a jej dopełnieniem ortogonalnym jest prosta rozpięta przez
wektor położenia $x$, o $\langle x,x\rangle_L=-1$.
Ten sam dowód zgodności i braku torsji co w
Stwierdzeniu~\ref{prop:koneksja-projekcja} daje
\[
 \nabla_XY=D_XY+\langle D_XY,x\rangle_Lx
          =D_XY-g(X,Y)x.
\]
Dla $\gamma(t)=\cosh(st)p+\sinh(st)v/s$ mamy
$\ddot\gamma=s^2\gamma$ i $g(\dot\gamma,\dot\gamma)=s^2$,
więc rzeczywiście $D\dot\gamma/\dd t=0$.
Wzór działa dla każdego $t$, co dowodzi zupełności geodezyjnej.

Uzasadnijmy dziedzinę odwrotnej funkcji hiperbolicznej.
Jeśli $p=(\sqrt{1+|z|^2},z)$ i $q=(\sqrt{1+|w|^2},w)$,
to nierówność Cauchy'ego--Schwarza daje
\[
 -\langle p,q\rangle_L
 \geq\sqrt{1+|z|^2}\sqrt{1+|w|^2}-|z|\,|w|\geq1.
\]
Druga nierówność po podniesieniu do kwadratu sprowadza się
do $(|z|-|w|)^2\geq0$; równość w obu zachodzi dokładnie
dla $z=w$. Dla dowolnego $q\in\mathcal H^n$ połóżmy
$s=\operatorname{arcosh}(-\langle p,q\rangle_L)$.
Gdy $q\ne p$, wektor
$v=\frac{s}{\sinh s}(q-\cosh(s)p)$
jest styczny w $p$, ma normę $s$ i spełnia
$\exp_p(v)=q$. Zatem $\exp_p$ jest także
surjekcją na $\mathcal H^n$.
Późniejszy Wniosek~\ref{cor:hopf-rinow-minimalna}
zapewni istnienie globalnej geodezyjnej minimalnej.
Jeśli $q=\exp_p(v)$, iloczyn z $p$ daje
$-\langle p,q\rangle_L=\cosh|v|_g$;
jednoznaczność dodatniego $|v|_g$
i wyrażenia \eqref{eq:exp-hiperboloida} pokazuje,
że $\exp_p$ jest różnowartościowe. Dlatego
\begin{equation}\label{eq:odleglosc-hiperboloida}
 d_{\mathcal H}(p,q)
 =\operatorname{arcosh}(-\langle p,q\rangle_L).
\end{equation}

W \emph{modelu kuli Poincarégo} z
\eqref{eq:hip-kula} skorzystajmy z izometrii
$P:B^n\to\mathcal H^n$ zapisanej w
Przykładzie~\ref{prz:trzy-modele-hiperboliczne}.
Ponieważ $P(0)=(1,0)$ i
$P^{-1}(t,z)=z/(t+1)$, obraz geodezyjnej
$\tau\mapsto(\cosh\tau,\sinh\tau\,e)$,
$|e|=1$, jest równy
\begin{equation}\label{eq:geodezyjna-dysk-radialna}
 \gamma_e(\tau)=\tanh(\tau/2)e.
\end{equation}
Wszystkie średnice dysku są więc geodezyjnymi
\emph{jako zbiory punktów}; parametr euklidesowy
na średnicy nie jest parametrem długości hiperbolicznej.

Pozostałe geodezyjne w dwuwymiarowym dysku są
łukami okręgów prostopadłych do brzegu $|u|=1$.
Oto rachunek wyjaśniający ten obraz. Geodezyjna
hiperboloidy leży w płaszczyźnie liniowej, którą
w $\R^{1,2}$ można zapisać jako
$\alpha t+b\cdot z=0$. Po podstawieniu
$P(u)=((1+|u|^2)/(1-|u|^2),2u/(1-|u|^2))$
jej obraz odwrotny spełnia
\[
 \alpha(1+|u|^2)+2b\cdot u=0.
\]
Jeśli $\alpha=0$, to średnica $b\cdot u=0$.
Jeśli $\alpha\ne0$, po oznaczeniu $c=-b/\alpha$
dostajemy
\begin{equation}\label{eq:geodezyjna-dysk-okrag}
 |u-c|^2=|c|^2-1.
\end{equation}
Środek $c$ leży poza dyskiem: płaszczyzna geodezyjnej
przecina hiperboloidę, więc równanie ma rozwiązanie
$0<|u|<1$; stąd
$|b|\geq |\alpha|(1+|u|^2)/(2|u|)>|\alpha|$.
Kwadrat promienia
okręgu wynosi $|c|^2-1$. Warunek ten oznacza
prostopadłe przecięcie z okręgiem jednostkowym:
w punkcie przecięcia $u$ równanie okręgu daje
$u\cdot(u-c)=0$, więc promienie obu okręgów
są prostopadłe. W wymiarze $n>2$ geodezyjna
leży w odpowiednim dwuwymiarowym przekroju
kuli i ten sam opis działa w tym przekroju.

Izometria $F:B^n\to\mathbb H^n$ z
Przykładu~\ref{prz:trzy-modele-hiperboliczne}
przenosi opis dyskowy na pionowe proste
i półokręgi w półprzestrzeni. Izometrie
zachowują pochodną kowariantną, więc obrazy
geodezyjnych są geodezyjnymi. Zachowują też
długości, skąd z~\eqref{eq:odleglosc-hiperboloida}
otrzymujemy użyteczne wzory odległości:
\begin{align}\label{eq:odleglosci-modele-hiperboliczne}
 \cosh d_B(u,v)
  &=1+\frac{2|u-v|^2}
      {(1-|u|^2)(1-|v|^2)},\\
 \cosh d_{\mathbb H}((x,y),(x',y'))
  &=1+\frac{|x-x'|^2+(y-y')^2}{2yy'}.\notag
\end{align}
Pierwszy wzór pochodzi z wstawienia $P(u),P(v)$
do iloczynu Minkowskiego; drugi z podstawienia
$F^{-1}(x,y),F^{-1}(x',y')$ do pierwszego.
W szczególności
$d_B(0,u)=2\operatorname{artanh}|u|	o\infty$
przy $|u|\uparrow1$: widoczny brzeg dysku,
podobnie jak prosta $y=0$, nie leży w przestrzeni
hiperbolicznej.
\end{przyklad}

\begin{figure}[H]
\centering
\begin{tikzpicture}[>=Latex,scale=.91]
  \fill[manifoldblue!5] (-3.5,0) rectangle (3.5,2.65);
  \draw[manifoldblue!70!black,very thick] (-3.5,0)--(3.5,0)
       node[right] {$y=0$};
  \draw[manifoldblue!45!black,->] (-3.25,0)--(-3.25,2.55)
       node[above] {$y$};
  \draw[manifoldred,very thick,->] (-1.60,.30)--(-1.60,2.23);
  \fill[manifoldred] (-1.60,.55) circle (2.5pt)
       node[left] {$(x_0,y_0)$};
  \fill[manifoldred] (-1.60,1.85) circle (2.5pt)
       node[left] {$(x_0,y_1)$};
  \draw[manifoldgreen!75!black,very thick,->]
       ({1.1+1.7*cos(170)},{1.7*sin(170)})
       arc[start angle=170,end angle=19,radius=1.7];
  \draw[manifoldgold!70!black,densely dashed]
       (-.6,0) arc[start angle=180,end angle=0,radius=1.7];
  \node[manifoldgreen!60!black] at (1.28,2.38)
       {półokrąg prostopadły do brzegu};
  \node[manifoldblue!70!black] at (1.55,-.47)
       {brzeg w nieskończonej odległości};
\end{tikzpicture}
\caption{Geodezyjne w modelu górnej półpłaszczyzny: pionowa półprosta i półokrąg przecinający brzeg $y=0$ pod kątem prostym. Brzeg nie należy do przestrzeni hiperbolicznej.}
\label{fig:geodezyjne-hiperboliczne}
\end{figure}

\begin{figure}[H]
\centering
\begin{tikzpicture}[>=Latex,scale=1]
  \begin{scope}[shift={(-2.75,0)}]
    \fill[manifoldblue!7] (0,0) circle (1.27);
    \draw[manifoldblue!70!black,very thick] (0,0) circle (1.27);
    \draw[manifoldgreen!75!black,very thick,->]
       (-1.04,0)--(1.06,0);
    \draw[manifoldred,very thick,->]
       ({2*1.27+sqrt(3)*1.27*cos(152)},
        {sqrt(3)*1.27*sin(152)})
       arc[start angle=152,end angle=208,
           radius={sqrt(3)*1.27}];
    \fill[manifoldred] ({2*1.27+sqrt(3)*1.27*cos(158)},
        {sqrt(3)*1.27*sin(158)}) circle (2.2pt)
        node[above left] {$p$};
    \fill[manifoldred] ({2*1.27+sqrt(3)*1.27*cos(202)},
        {sqrt(3)*1.27*sin(202)}) circle (2.2pt)
        node[below left] {$q$};
    \node[manifoldblue!75!black] at (0,-1.82)
       {dysk Poincarégo};
  \end{scope}
  \begin{scope}[shift={(2.65,0)}]
    % Dokładny przekrój z_2=0: t=cosh(u), z_1=sinh(u).
    % Rysunek ma skalę poziomą 0.6 oraz przesunięcie pionowe -1.2.
    \fill[manifoldgold!9] (-1.75,-1.35) rectangle (1.85,1.65);
    \draw[manifoldblue!55!black,->] (-1.6,-1.2)--(1.6,-1.2)
      node[right] {$z_1$};
    \draw[manifoldblue!55!black,->] (0,-1.4)--(0,1.60)
      node[above] {$t$};
    \draw[manifoldblue!60!black,thick]
      plot[domain=-1.6:1.6,samples=90,variable=\u]
        ({.6*sinh(\u)},{cosh(\u)-1.2});
    \draw[manifoldred,very thick,->]
      plot[domain=-.65:1.35,samples=65,variable=\u]
        ({.6*sinh(\u)},{cosh(\u)-1.2});
    \fill[manifoldred] (0,-.2) circle (2.2pt)
      node[below right] {$p_0$};
    \fill[manifoldred] ({.6*sinh(1)},{cosh(1)-1.2}) circle (2.2pt)
      node[right] {$q_0$};
    \fill[manifoldgold!80!black] (0,-1.2) circle (1.8pt)
      node[below left] {$0$};
    \node[manifoldblue!75!black] at (0,-1.82)
      {przekrój hiperboloidy: $z_2=0$};
  \end{scope}
\end{tikzpicture}
\caption{Po lewej geodezyjne dysku: średnica i łuk okręgu
prostopadłego do brzegu (końce na brzegu nie należą do geodezyjnej).
Po prawej dokładny przekrój hiperboloidy $t^2-z_1^2-z_2^2=1$,
$t>0$, płaszczyzną liniową $z_2=0$. Czerwona geodezyjna ma
postać $(\cosh u,\sinh u,0)$; zaznaczono $p_0$ dla $u=0$
i $q_0$ dla $u=1$. Płaszczyzna przekroju zawiera początek $0$.}
\label{fig:geodezyjne-dysk-hiperboloida}
\end{figure}

\section{Twierdzenie Hopfa--Rinowa}
\label{sec:hopf-rinow}

W lokalnym twierdzeniu minimalnym odcinek był
krótki. Żeby mieć geodezyjną minimalną między \emph{dowolnymi}
punktami, potrzebujemy warunku, który zabrania
geodezyjnej uciec z rozmaitości w skończonym
czasie. Dla spójnej rozmaitości riemannowskiej
równoważne są dwa pozornie różne rodzaje
zupełności.

\begin{definicja}[Dwa znaczenia zupełności]
\index{dwa znaczenia zupelnosci@Dwa znaczenia zupełności}
\label{def:zupelnosc-geodezyjna}
$(M,g)$ jest \emph{metrycznie zupełna}, gdy
każdy ciąg Cauchy'ego dla $d_g$ ma granicę w $M$.
Jest \emph{geodezyjnie zupełna}, gdy każdą
geodezyjną można przedłużyć do parametru
z całej $\R$.
\end{definicja}

Do dowodu potrzebne jest spostrzeżenie o narożniku.
Jeżeli dwie jednostkowe geodezyjne składają się w drogę
minimalną, ich wektory styczne w punkcie sklejenia $x$
muszą być równe. Oznaczmy je przez $T_-,T_+$,
a chwilę sklejenia przez $t_0$.
W mapie wokół $x$ wybieramy ciągłe, kawałkami gładkie
pole wariacyjne o małym nośniku wokół chwili sklejenia,
takie że $V(t_0)=T_+-T_-$. Mnożenie jego współrzędnych
przez funkcję odcinającą realizuje taką wariację przez
dodanie $sV$ we współrzędnych, przy nieruchomych końcach.
Dla małego $s$ oba kawałki pozostają regularne.
Wzór pierwszej wariacji na obu geodezyjnych ma zerowe
składniki całkowe, a wyraz narożnikowy wynosi
\[
 g(T_+-T_-,T_--T_+)=-|T_+-T_-|_g^2.
\]
Jeśli styczne są różne, długość maleje dla małego $s>0$,
co przeczy minimalności. Przy jednakowych niezerowych
prędkościach ten argument stosujemy po wspólnym skalowaniu.

\begin{lemat}[Jeden punkt z pełnym odwzorowaniem wykładniczym]
\label{lem:exp-p-minimalna-kazdy-punkt}
Niech $M$ będzie spójną rozmaitością riemannowską bez brzegu.
Jeżeli dla pewnego $p\in M$ odwzorowanie $\exp_p$ jest określone na całym $T_pM$,
to każdy $q\in M$ można połączyć z $p$
geodezyjną o długości $d_g(p,q)$.
\end{lemat}
\begin{proof}
Jeśli $q=p$, bierzemy stałą geodezyjną.
Niech $r=d_g(p,q)>0$. Wybierzmy małą
kulę $B^g_\varepsilon(p)$ spełniającą
\eqref{eq:male-kule-normalne} i
$0<\delta<\min(\varepsilon,r)$.
Sfera $S_\delta(p)=\exp_p\{v:|v|=\delta\}$
jest zwarta jako obraz zwykłej zwartej sfery.
Funkcja $z\mapsto d_g(z,q)$ jest ciągła,
więc osiąga na niej minimum w punkcie $m$.
Każda droga z $p$ do $q$ przecina tę sferę.
Jej część do pierwszego przecięcia ma długość
co najmniej $\delta$, pozostała zaś co najmniej
$d_g(z,q)\geq d_g(m,q)$. Z drugiej strony
nierówność trójkąta daje
\begin{equation}\label{eq:hopf-pierwsza-sfera}
 r=\delta+d_g(m,q).
\end{equation}
Niech $\gamma:[0,\infty)\to M$ będzie
jednostkową geodezyjną od $p$ przez $m$.
Istnieje dla wszystkich $t\geq0$: dla danego
$t$ punkt $\gamma(t)=\exp_p(t\dot\gamma(0))$
jest określony, a jednoznaczność skleja te odcinki.
Rozważmy zamknięty zbiór
\[
 A=\{t\in[0,r]:t+d_g(\gamma(t),q)=r\}.
\]
Z \eqref{eq:hopf-pierwsza-sfera} wiemy, że
$\delta\in A$. Jeśli największy punkt
$t_*$ zbioru $A$ byłby mniejszy od $r$,
połóżmy $x=\gamma(t_*)$ i wybierzmy małą
normalną sferę $S_\eta(x)$ o promieniu
$0<\eta<d_g(x,q)$, tak by $\eta$ mieściło się
również w normalnym otoczeniu $x$.
Niech $m'$ minimalizuje $z\mapsto d_g(z,q)$
na $S_\eta(x)$. To samo rozumowanie
z pierwszym wyjściem z kuli daje
\begin{equation}\label{eq:hopf-druga-sfera}
 d_g(x,q)=\eta+d_g(m',q),\qquad
 r=t_*+\eta+d_g(m',q).
\end{equation}

Nierówność trójkąta i droga przez $x,m'$ dają
\[
 r\leq d_g(p,m')+d_g(m',q)
   \leq t_*+\eta+d_g(m',q)=r.
\]
Zachodzi więc równość: sklejenie $\gamma|_{[0,t_*]}$
z krótką radialną geodezyjną z $x$ do $m'$
minimalizuje długość od $p$ do $m'$.
Obie części są jednostkowe, zatem spostrzeżenie
o narożniku wymusza zgodność ich stycznych w $x$.
Jednoznaczność równania geodezyjnej daje
$m'=\gamma(t_*+\eta)$. Wtedy
\eqref{eq:hopf-druga-sfera} mówi, że
$t_*+\eta\in A$, sprzeczność z maksymalnością
$t_*$. Musi więc być $r\in A$,
czyli $d_g(\gamma(r),q)=0$ i $\gamma(r)=q$.
Ponieważ $\gamma$ ma prędkość $1$,
jej długość na $[0,r]$ wynosi $r=d_g(p,q)$.
\end{proof}

\begin{twierdzenie}[Hopfa--Rinowa]
\label{thm:hopf-rinow-geodezyjne}
Niech $(M,g)$ będzie niepustą, spójną rozmaitością
riemannowską bez brzegu. Równoważne są warunki:
\begin{enumerate}[label=(\roman*)]
\item $(M,d_g)$ jest zupełną przestrzenią metryczną;
\item każda maksymalna geodezyjna jest określona
      na całej $\R$;
\item dla pewnego $p\in M$ mamy $D_p=T_pM$,
      czyli $\exp_p$ jest określone wszędzie;
\item każdy domknięty i ograniczony w metryce
      $d_g$ podzbiór $M$ jest zwarty.
\end{enumerate}
Gdy warunki zachodzą, w (iii) można napisać
\emph{dla każdego} $p\in M$.
\end{twierdzenie}
\begin{proof}
Pokażemy kolejno (i)$\Rightarrow$(ii)
$\Rightarrow$(iii)$\Rightarrow$(iv)$\Rightarrow$(i).

\emph{(i)$\Rightarrow$(ii).}
Jeśli jednostkowa geodezyjna
$\gamma:[0,b)\to M$ miałaby maksymalny
skończony prawy koniec $b$, to
$d_g(\gamma(s),\gamma(t))\leq|s-t|$.
Dla dowolnego ciągu $t_j\uparrow b$ ciąg $\gamma(t_j)$
jest Cauchy'ego, więc zbiega do pewnego $z\in M$.
Oszacowanie odległości przez $|s-t|$ daje wtedy
$\gamma(t)\to z$ dla wszystkich $t\uparrow b$.
W mapie wokół $z$ macierz $G$ jest dodatnio
określona i na mniejszym zwartym otoczeniu
jednostajnie porównywalna z macierzą
euklidesową. Stałość prędkości geodezyjnej
ogranicza więc jej współczynniki $\dot x^i$.
Symbole $\Gamma^k_{ij}$ są tam ograniczone;
równanie~\eqref{eq:geodezyjna-ode} ogranicza
również $\ddot x^k$. Stąd $x(t)$ i $\dot x(t)$
mają granice przy $t\uparrow b$.
Lokalne twierdzenie ODE, uruchomione od tych
danych w chwili $b$, przedłuża geodezyjną
za $b$ -- sprzeczność. Analogicznie na lewym
końcu. Dla dowolnej stałej prędkości
przeskalowujemy parametr; geodezyjne o zerowej prędkości
są stałe i istnieją przez cały czas.

\emph{(ii)$\Rightarrow$(iii).}
Każda $\gamma_{p,v}$ istnieje w chwili $1$,
więc $D_p=T_pM$ dla każdego $p$.

\emph{(iii)$\Rightarrow$(iv).}
Lemat~\ref{lem:exp-p-minimalna-kazdy-punkt}
daje dla każdego $q$ geodezyjną od $p$
o długości $d_g(p,q)$. Jeśli $C$ jest domknięte
i ograniczone, wybierzmy $R$ z
$d_g(p,q)\leq R$ dla wszystkich $q\in C$.
Każdy $q\in C$ ma postać $\exp_p(v_q)$
z $|v_q|\leq R$. Zatem
$C\subset\exp_p(\overline B_R(0))$.
Domknięta kula w skończenie wymiarowej
przestrzeni $T_pM$ jest zwarta, jej obraz
też jest zwarty, a domknięty podzbiór
tego obrazu jest zwarty. To dowodzi (iv).

\emph{(iv)$\Rightarrow$(i).}
Ciąg Cauchy'ego jest ograniczony;
domknięcie zbioru jego wyrazów jest
domknięte i ograniczone, więc zwarte.
Istnieje zbieżny podciąg, a własność Cauchy'ego
zmusza cały ciąg do tej samej granicy.
\end{proof}

\begin{wniosek}[Globalna geodezyjna minimalna]
\label{cor:hopf-rinow-minimalna}
Jeżeli $(M,g)$ jest spójna i zupełna, to
dowolne $p,q\in M$ łączy geodezyjna
o długości $d_g(p,q)$. W szczególności
każda zwarta rozmaitość riemannowska
jest geodezyjnie zupełna.
\end{wniosek}
\begin{proof}
Pierwszą tezę daje (ii)$\Rightarrow$(iii)
oraz Lemat~\ref{lem:exp-p-minimalna-kazdy-punkt}.
Dla zwartej $M$ argument stosujemy w każdej składowej
spójnej, która jest domknięta i zwarta. Każdy ciąg w niej
ma zbieżny podciąg; ciąg Cauchy'ego zbiega
do granicy tego podciągu. Zatem (i) zachodzi
i stosujemy Twierdzenie~\ref{thm:hopf-rinow-geodezyjne}.
\end{proof}

\begin{przyklad}[Dlaczego potrzebujemy zupełności]
W $M=\R^2\setminus\{0\}$ z metryką euklidesową
geodezyjna $\gamma(t)=(1-t,0)$ dociera do
usuniętego punktu w skończonym czasie $t=1$.
Nie można jej przedłużyć jako krzywej w $M$.
Punkty $(1/n,0)$ tworzą ciąg Cauchy'ego bez
granicy w $M$. Między $(-1,0)$ i $(1,0)$
infimum długości wynosi $2$, ale żadna krzywa
w $M$ go nie osiąga. Dolne ograniczenie $2$ daje
nierówność euklidesowa. Dwa odcinki osi, od $-1$ do
$-\varepsilon$ i od $\varepsilon$ do $1$, połączone
półokręgiem promienia $\varepsilon$, mają łączną długość
$2-2\varepsilon+\pi\varepsilon\to2$.
Równość w nierówności euklidesowej wymagałaby przebiegania
odcinka od $(-1,0)$ do $(1,0)$ bez cofania, a taki odcinek
przechodzi przez usunięty punkt.
\end{przyklad}

\begin{figure}[htbp]
\centering
\begin{tikzpicture}[>=Latex,scale=1]
  \fill[manifoldblue!6,rounded corners=8pt]
       (-3,-1.10) rectangle (3,1.10);
  \draw[manifoldblue!50,rounded corners=8pt]
       (-3,-1.10) rectangle (3,1.10);
  \coordinate (a) at (-2.3,0);
  \coordinate (b) at (2.3,0);
  \draw[manifoldred!75!black,very thick,->] (1.78,0)--(.30,0);
  \draw[manifoldred!75!black,very thick,->] (-1.78,0)--(-.30,0);
  \draw[manifoldgold!85!black,thick,densely dashed] (a)--(b);
  \draw[manifoldblue!70!black,thick]
       (a) .. controls (-1.05,.76) and (1.05,.76) .. (b);
  \draw[manifoldblue!70!black,thick]
       (a) .. controls (-1.05,-.76) and (1.05,-.76) .. (b);
  \fill[white] (0,0) circle (5.2pt);
  \draw[manifoldred!85!black,thick] (0,0) circle (5.2pt);
  \node[manifoldred!85!black] at (0,.88) {$0\notin M$};
  \fill[manifoldblue!80!black] (a) circle (2pt) node[below] {$(-1,0)$};
  \fill[manifoldblue!80!black] (b) circle (2pt) node[below] {$(1,0)$};
\end{tikzpicture}
\caption{Na przebitej płaszczyźnie prosta geodezyjna może dotrzeć do usuniętego punktu w skończonym czasie. Drogi łączące punkty po przeciwnych stronach otworu omijają go; ich długości mogą zbliżać się do $2$, lecz żadna nie osiąga tej wartości.}
\label{fig:geodezyjna-przebita-plaszczyzna}
\end{figure}

\section{Schwarzschild i Kerr: geodezyjne czasoprzestrzeni}
\label{sec:geodezyjne-schwarzschild-kerr}

Obie metryki z \eqref{eq:schwarzschild} i
\eqref{eq:kerr} są \emph{lorentzowskie}.
Równanie~\eqref{eq:geodezyjna-ode} i $\exp_p$
nadal mają sens lokalnie, podobnie jak ich
dowody z ODE. Nie ma jednak dodatniej normy
$\sqrt{g(v,v)}$ dla wszystkich wektorów,
więc odległość $d_g$ z Rozdziału~\ref{ch:metryka-riemannowska},
lemat o minimalności długości i twierdzenie
Hopfa--Rinowa nie stosują się do \emph{całej}
czasoprzestrzeni.

Dla czasopodobnej krzywej z parametrem $\lambda$
jej \emph{czas własny} jest równy
\begin{equation}\label{eq:czas-wlasny}
 \tau(\gamma)=\int\sqrt{-g(\dot\gamma,\dot\gamma)}
                    \,\dd\lambda.
\end{equation}
Przy sygnaturze $(-,+,+,+)$ czasopodobna
geodezyjna między dostatecznie bliskimi
punktami \emph{maksymalizuje} czas własny
wśród przyszłościowych krzywych czasopodobnych o tych samych
końcach, zawartych w odpowiednio małym otoczeniu normalnym. W płaskim modelu
$g=-\dd t^2+|\dd x|^2$, po przejściu do
układu spoczynkowego końców, dla
$x(a)=x(b)$ mamy wprost
\[
 \tau(\alpha)=\int_a^b
 \sqrt{\dot t^2-|\dot x|^2}\,\dd\lambda
 \leq\int_a^b\dot t\,\dd\lambda=t(b)-t(a),
\]
z równością dla prostej o stałym $x$.
W zakrzywionej czasoprzestrzeni ten argument
ma lokalną wersję. Dowód lematu
Gaussa~\ref{lem:gaussa-exp} używa jedynie
zgodności z metryką i braku torsji, więc
działa również dla metryki lorentzowskiej.
Wybieramy orientację czasu w małym otoczeniu normalnym $U$
punktu $p$. Możemy je zmniejszyć tak, aby istniała współrzędna
$x^0$ z $\dd x^0$ dodatnią na wszystkich niezerowych
przyszłościowych wektorach przyczynowych. Wystarczy wybrać taką współrzędną
w $p$ i zachować czasopodobność jej gradientu przez ciągłość.

Dlaczego przyszłościowa krzywa czasopodobna wychodząca z $p$
i pozostająca w $U$ leży w obrazie przyszłego stożka?
Niech $Q(\exp_p v)=g_p(v,v)$ i $R_{\exp_p v}=\dd(\exp_p)_v(v)$.
Lemat Gaussa daje $\dd Q(W)=2g(R,W)$.
Początkowo krzywa leży w przyszłym stożku, gdyż
$\dd(\exp_p)_0=\id$ i jej początkowa prędkość jest czasopodobna.
Gdyby po raz pierwszy doszła do niezerowego punktu jego
świetlnego brzegu, mielibyśmy tam $Q=0$ i $Q'\geq0$.
Jednak $R$ jest wtedy niezerowym przyszłościowym wektorem
światłopodobnym, zatem $Q'=2g(R,\dot\alpha)<0$.
Ostatnia nierówność wynika w bazie lorentzowskiej z
$R=(r,\xi)$, $r=|\xi|>0$, i $\dot\alpha=(a,\eta)$,
$a>|\eta|$: $-ra+\xi\cdot\eta<0$.
Powrót do wierzchołka $p$ wyklucza ścisły wzrost $x^0$.
Krzywa pozostaje więc w przyszłym stożku.

We wnętrzu tego stożka
piszemy $v=\tau u$, gdzie $g_p(u,u)=-1$,
$\tau>0$. Ortogonalne do czasopodobnego
kierunku $u$ wektory mają dodatni kwadrat
normy. Lemat Gaussa daje więc w tych
współrzędnych schemat metryki
$g=-\dd\tau^2+h_{ab}(\tau,u)\dd u^a\dd u^b$
z dodatnio określonym $h_{ab}$.
Radialne pole $\partial_\tau$ jest przyszłościowe,
jednostkowe i czasopodobne. Dla przyszłościowej krzywej
czasopodobnej mamy więc $\dot\tau=-g(\partial_\tau,\dot\alpha)>0$ i
\[
 \sqrt{-g(\dot\alpha,\dot\alpha)}
 =\sqrt{(\dot\tau)^2-h_{ab}\dot u^a\dot u^b}
 \leq\dot\tau.
\]
Całkujemy od dodatniego parametru i przechodzimy do początku
krzywej, gdzie $\tau\to0$. Dostajemy $\tau(\alpha)
\leq\tau(q)$; radialna geodezyjna osiąga
równość. Tym samym lokalna maksymalność
dotyczy dróg mieszczących się w takim
otoczeniu przy ustalonych końcach.
Nie jest to twierdzenie o \emph{minimalizowaniu}
riemannowskiej odległości. Dla geodezyjnych
światłopodobnych czas własny wynosi zero.

\begin{przyklad}[Ruch w metryce Schwarzschilda]
\label{prz:geodezyjne-schwarzschild}
Pracujemy w obszarze $r>2M$, $M>0$. Z uwagi na symetrię
sferyczną możemy wybrać płaszczyznę ruchu $\theta=\pi/2$:
obracamy początkowe dane przestrzenne do jednej płaszczyzny
przez środek. Odbicie względem niej jest izometrią
zachowującą te dane, więc jednoznaczność geodezyjnej
utrzymuje ruch w tej płaszczyźnie.
Na geodezyjnej z parametrem afinicznym
Lagrangian $\mathcal L=\tfrac12g_{\mu\nu}
\dot x^\mu\dot x^\nu$ nie zależy od $t,\phi$.
Równania Eulera--Lagrange'a dają więc
stałe wielkości
\begin{equation}\label{eq:stale-schwarzschild}
 E=f(r)\dot t,\qquad L=r^2\dot\phi,
 \qquad f(r)=1-\frac{2M}{r}.
\end{equation}
Niech $g(\dot\gamma,\dot\gamma)=-\kappa$:
$\kappa=1$ dla czasopodobnej geodezyjnej
parametryzowanej czasem własnym i $\kappa=0$
dla światłopodobnej. Podstawienie
$\dot t=E/f$, $\dot\phi=L/r^2$ do metryki daje
\begin{equation}\label{eq:radialne-schwarzschild}
 -\kappa=-\frac{E^2}{f}+\frac{\dot r^2}{f}
                   +\frac{L^2}{r^2}
 \quad\Longleftrightarrow\quad
 \dot r^2+f\left(\kappa+\frac{L^2}{r^2}\right)=E^2.
\end{equation}
To pierwsza całka równania radialnego. W punktach zwrotnych
nie wolno dzielić jej pochodnej przez $\dot r=0$.
Pełne równanie geodezyjnej, po wstawieniu symboli
z \eqref{eq:gamma-schwarzschild} i powyższej pierwszej całki, daje
\[
 \ddot r=-\tfrac12 f'\left(\kappa+\frac{L^2}{r^2}\right)
                  +\frac{fL^2}{r^3}
        =-\tfrac12\frac{\dd}{\dd r}
                    \left[f\left(\kappa+\frac{L^2}{r^2}\right)\right].
\]
Ten wzór obowiązuje także przy $\dot r=0$.
Przy $\kappa=0$, $L\ne0$ stały promień
orbity światłopodobnej wymaga $\dot r=0$
i pochodnej potencjału
$V(r)=L^2(1-2M/r)/r^2$ równej zeru.
Ponieważ
$V'(r)=2L^2(3M-r)/r^4$, dostajemy
$r=3M$ i z \eqref{eq:radialne-schwarzschild}
$E^2=L^2/(27M^2)$. To przykład geodezyjnej,
która może krążyć po sferze fotonowej;
nie jest to okrąg minimalizujący długość
w sensie riemannowskim.
\end{przyklad}

\begin{figure}[htbp]
\centering
\begin{tikzpicture}[>=Latex,scale=1]
  \shade[inner color=manifoldgold!50,outer color=manifoldgold!8]
       (0,0) circle (2.10);
  \fill[manifoldblue!75!black] (0,0) circle (.72);
  \draw[manifoldblue!35!black,thick] (0,0) circle (.72);
  \draw[manifoldgold!80!black,densely dashed]
       (0,0) circle (1.08);
  \draw[manifoldred!85!black,very thick,->]
       (0:1.08) arc[start angle=0,end angle=116,radius=1.08];
  \draw[manifoldred!85!black,very thick,->]
       (145:1.08) arc[start angle=145,end angle=330,radius=1.08];
  \fill[manifoldred] (1.08,0) circle (2.5pt)
       node[right] {$\gamma$};
  \draw[manifoldblue!60!black,->]
       (2.30,-1.0)--(-25:1.08)
       node[pos=0,below] {$r=3M$};
  \node[white] at (0,0) {$r\leq2M$};
\end{tikzpicture}
\caption{Schemat rzutu na współrzędne $(r,\phi)$ równikowej geodezyjnej światłopodobnej Schwarzschilda o stałym $r=3M$. Ciemny dysk oznacza $r\leq2M$ w tym rysunku współrzędnych; okrąg $r=3M$ nie przedstawia geodezyjnej metryki indukowanej na przekroju przestrzennym.}
\label{fig:geodezyjna-schwarzschild-foton}
\end{figure}

\begin{przyklad}[Ruch równikowy w metryce Kerra]
\label{prz:geodezyjne-kerr}
Przyjmujemy $M>0$, $|a|\leq M$ i $r>r_+=M+\sqrt{M^2-a^2}$.
Płaszczyzna $\theta=\pi/2$ jest niezmiennicza
przez symetrię odbicia $\theta\mapsto\pi-\theta$:
geodezyjna startująca na niej ze
$\dot\theta=0$ na niej pozostaje.
Niech $E=-p_t$, $L=p_\phi$, gdzie
$p_\mu=g_{\mu\nu}\dot x^\nu$; są to stałe,
bo metryka nie zależy od $t,\phi$.
W tej płaszczyźnie $\Sigma=r^2$. Oznaczmy
\[
 A=\dot t-a\dot\phi,
 \qquad B=(r^2+a^2)\dot\phi-a\dot t.
\]
Z postaci \eqref{eq:kerr} przez różniczkowanie
Lagrangianu względem $\dot t,\dot\phi$
otrzymujemy
\[
 E=\frac{\Delta A+aB}{r^2},
 \qquad
 L=\frac{a\Delta A+(r^2+a^2)B}{r^2}.
\]
Odejmujemy $aE$ od $L$: dostajemy
$B=L-aE$, a stąd $\Delta A=P(r)$,
gdzie $P(r)=E(r^2+a^2)-aL$.
Warunek $g(\dot\gamma,\dot\gamma)=-\kappa$
daje po pomnożeniu przez $r^2\Delta$
\begin{equation}\label{eq:radialne-kerr}
 r^4\dot r^2=P(r)^2
        -\Delta\bigl(\kappa r^2+(L-aE)^2\bigr),
 \qquad \kappa\in\{0,1\}.
\end{equation}
Możemy również odzyskać pozostałe pochodne:
\begin{equation}\label{eq:tphi-kerr}
 r^2\dot t=\frac{r^2+a^2}{\Delta}P+a(L-aE),
 \qquad
 r^2\dot\phi=\frac a\Delta P+(L-aE).
\end{equation}
Wzory dotyczą geodezyjnych równikowych;
poza tą płaszczyzną ruch wymaga jeszcze
jednej stałej związanej z kątem $\theta$.
Przy $a=0$ równanie~\eqref{eq:radialne-kerr}
redukuje się do~\eqref{eq:radialne-schwarzschild}.
Także tutaj pierwsza całka radialna nie zastępuje pełnego
równania geodezyjnej przy poszukiwaniu rozwiązań stałych
lub przedłużaniu przez punkty zwrotne.

Dla kołowej geodezyjnej równikowej o
$\Omega=\dot\phi/\dot t$ i $\dot r=0$
radialne równanie geodezyjnej daje
\[
 \partial_rg_{tt}+2\Omega\partial_rg_{t\phi}
                 +\Omega^2\partial_rg_{\phi\phi}=0.
\]
Na równiku
$g_{tt}=-1+2M/r$,
$g_{t\phi}=-2Ma/r$,
$g_{\phi\phi}=r^2+a^2+2Ma^2/r$;
po pomnożeniu równania przez $r^2/2$
otrzymujemy
\[
 -M+2Ma\Omega+(r^3-Ma^2)\Omega^2=0,
 \qquad
 \Omega_\pm=\frac{\pm\sqrt M}{r^{3/2}\pm a\sqrt M}.
\]
To \emph{kandydaci} na orbity kołowe: dla
czasopodobnej trzeba ponadto sprawdzić
$g(\partial_t+\Omega_\pm\partial_\phi,
   \partial_t+\Omega_\pm\partial_\phi)<0$
w rozważanym promieniu. Przy $a=0$
dostajemy $\Omega_\pm=\pm\sqrt{M/r^3}$.
\end{przyklad}

\begin{figure}[htbp]
\centering
\begin{tikzpicture}[>=Latex,scale=1]
  \shade[inner color=manifoldblue!24,outer color=white]
       (0,0) circle (2.16);
  \fill[manifoldblue!75!black] (0,0) circle (.64);
  \draw[manifoldblue!50!black,thick] (0,0) circle (.64);
  \draw[manifoldgold!75!black,thick,densely dashed]
       (0,0) circle (1.42);
  \draw[manifoldred,very thick,->]
       (30:1.42) arc[start angle=30,end angle=146,radius=1.42];
  \draw[manifoldgreen!75!black,very thick,->]
       (330:1.42) arc[start angle=330,end angle=214,radius=1.42];
  \node[manifoldred!85!black] at (-2.23,1.24)
       {$\Omega_+>0$};
  \node[manifoldgreen!70!black] at (2.30,-.75)
       {$\Omega_-<0$};
  \draw[manifoldgold!85!black,thick,->]
       (-.29,-.31) arc[start angle=225,end angle=520,radius=.43];
  \node[manifoldgold!80!black] at (.13,.15) {$a>0$};
  \node[manifoldblue!80!black] at (2.22,1.40)
       {$r=r_0$};
\end{tikzpicture}
\caption{Schemat rzutu na płaszczyznę równikową dwóch kierunków kołowej geodezyjnej Kerra przy dodatnim $a$ i odpowiednio dużym $r_0$. Strzałki $\Omega_+$ i $\Omega_-$ oznaczają ruch zgodny z obrotem i przeciwny. Ciemny obszar schematycznie zaznacza wnętrze $r_+$, a okrąg przerywany jest stałym $r_0$ we współrzędnych Boyera--Lindquista; nie jest to zanurzenie zakrzywionej czasoprzestrzeni.}
\label{fig:geodezyjna-kerr-kolowa}
\end{figure}

\subsection*{Riemannowskie przekroje przestrzenne: co w nich minimalizuje?}
\label{subsec:przekroje-przestrzenne-geodezyjne}

Można osobno wybrać przekrój $t=\mathrm{const}$
i jego \emph{indukowaną metrykę riemannowską}
$h$. To inny problem geometryczny niż
geodezyjne pełnej czasoprzestrzeni.
Rysunek~\ref{fig:geodezyjne-przestrzen-czas} zestawia
przestrzenną długość z czasem własnym; prawe okno
pokazuje płaską czasoprzestrzeń, aby oddzielić te
dwa rodzaje pomiaru przed obliczeniami metryk
Schwarzschilda i Kerra.

\begin{figure}[H]
\centering
\begin{tikzpicture}[>=Latex,scale=.95]
  \fill[manifoldblue!6] (-2.9,0) circle (1.55);
  \draw[manifoldblue!70!black,thick] (-2.9,0) circle (1.55);
  \fill[manifoldgold!48] (-2.9,0) circle (.46);
  \draw[manifoldgold!85!black,thick] (-2.9,0) circle (.46);
  \draw[manifoldblue,very thick] (-2.9,0) ++(.76,0)--++(.69,0);
  \fill[manifoldred] (-2.14,0) circle (2pt) node[below] {$p$};
  \fill[manifoldred] (-1.45,0) circle (2pt) node[below] {$q$};
  \draw[manifoldred,thick,densely dashed]
    (-2.14,0) .. controls (-2.3,.6) and (-1.52,.68) .. (-1.45,0);
  \node[manifoldgold!65!black] at (-2.9,.1) {$r_+$};
  \node[manifoldblue!80!black] at (-2.9,2.02) {Przekrój $t=\mathrm{const}$};
  \draw[manifoldblue!55!black,->] (2.1,-1.35)--(2.1,1.55) node[above] {$t$};
  \draw[manifoldblue!55!black,->] (.85,-1.1)--(3.45,-1.1) node[right] {$x$};
  \coordinate (lp) at (1.55,-.88);
  \coordinate (lq) at (2.65,1.13);
  \draw[manifoldblue!45,densely dashed] (lp)--++(-1.02,1.02);
  \draw[manifoldblue!45,densely dashed] (lp)--++(1.02,1.02);
  \draw[manifoldgreen!75!black,very thick] (lp)--(lq);
  \draw[manifoldred,thick] (lp)
    .. controls (1.48,.20) and (2.20,.52) .. (lq);
  \fill[manifoldblue!80!black] (lp) circle (2pt) node[below left] {$P$};
  \fill[manifoldblue!80!black] (lq) circle (2pt) node[above] {$Q$};
  \node[manifoldblue!80!black] at (2.1,2.02) {Czas własny $\tau$};
\end{tikzpicture}
\caption{Po lewej odcinek radialny minimalizuje długość w riemannowskim przekroju przestrzennym. Po prawej prosta czasopodobna maksymalizuje czas własny między ustalonymi zdarzeniami w płaskiej czasoprzestrzeni. Schemat nie przedstawia zanurzenia metryki Kerra.}
\label{fig:geodezyjne-przestrzen-czas}
\end{figure}

W zewnętrznym obszarze Schwarzschilda
\[
 h_{\mathrm{Sch}}=f^{-1}\dd r^2
    +r^2(\dd\theta^2+\sin^2\theta\,\dd\phi^2).
\]
Na równiku krzywa o stałych $\theta=\pi/2$
i $\phi$ oraz zmiennym $r$ ma długość
\begin{equation}\label{eq:radialna-przestrzenna-schwarz}
 \ell_{\mathrm{Sch}}(r_0,r_1)
  =\left|\int_{r_0}^{r_1}
               \frac{\dd r}{\sqrt{1-2M/r}}\right|.
\end{equation}
Dowolna droga w tym przekroju między punktami
o tych samych kątach ma długość co najmniej
$\int|\dot r|/\sqrt f\,\dd s\geq\ell_{\mathrm{Sch}}$;
monotoniczna krzywa radialna osiąga równość.

Dla Kerra na $t=\mathrm{const}$ i $\theta=\pi/2$
składowe metryki indukowanej wynoszą
\[
 h_{rr}=\frac{r^2}{\Delta},\qquad
 h_{\phi\phi}=r^2+a^2+\frac{2Ma^2}{r}.
\]
Stąd radialna długość ma wzór
\begin{equation}\label{eq:radialna-przestrzenna-kerr}
 \ell_{\mathrm K}(r_0,r_1)
   =\left|\int_{r_0}^{r_1}
                   \frac r{\sqrt\Delta}\,\dd r\right|.
\end{equation}
Ta droga też jest najkrótsza w \emph{tym przekroju}
między równikowymi punktami o tych samych
kątach: dla dowolnej drogi $\Sigma\geq r^2$
i $h_{rr}=\Sigma/\Delta$, zatem
$|\dot\alpha|_h\geq r|\dot r|/\sqrt\Delta$.
Po scałkowaniu dostajemy dolne ograniczenie
\eqref{eq:radialna-przestrzenna-kerr}, osiągane
przez krzywą radialną na równiku. Wobec
Wniosku~\ref{cor:minimalna-geodezyjna}
obie radialne krzywe, po parametryzacji
długością łuku, są geodezyjnymi swoich
\emph{przestrzennych} metryk $h$.

Te zewnętrzne przekroje $r>r_+$ nie są
automatycznie zupełne. Dla Schwarzschilda
całka~\eqref{eq:radialna-przestrzenna-schwarz}
do pominiętego brzegu $r=2M$ jest skończona.
Dla Kerra o $|a|<M$ mamy
$\Delta=(r-r_+)(r-r_-)$, $r_-<r_+$,
więc całka~\eqref{eq:radialna-przestrzenna-kerr}
do $r_+$ również jest skończona.
Wynika to z zachowania całkowanego wyrażenia
jak $C/\sqrt{r-r_+}$ blisko horyzontu.
Horyzont nie należy do opisywanego przekroju;
ciąg punktów na ustalonym promieniu równikowym,
z $r\downarrow r_+$, jest w nim ciągiem Cauchy'ego bez granicy. Nie przeczy
to Hopfowi--Rinowowi, ponieważ jego hipoteza
zupełności tutaj nie zachodzi.

\par\smallskip
{\footnotesize\noindent\emph{Dalsza lektura:}
\href{https://math.uchicago.edu/~may/REU2016/REUPapers/Spiegel.pdf}
{dowód Hopfa--Rinowa},
\href{https://www.math.toronto.edu/mein/teaching/LectureNotes/rieall.pdf}
{geodezyjne i współrzędne normalne},
\href{https://math.uchicago.edu/~dannyc/courses/kleinian_groups_2015/kleinian_groups.pdf}
{geodezyjne w modelach hiperbolicznych},
\href{https://ocw.mit.edu/courses/8-033-relativity-fall-2006/resources/lecture21_bh1/}
{orbity Schwarzschilda w notatkach MIT},
\href{https://sites.science.oregonstate.edu/coursewikis/GGR/book/content/kerr}
{metryka Kerra}.\par}
